%% ============================================================
%%  Reproducibility Package README
%%  Agricultural Loans and Expenditure Allocation in Nigeria:
%%  Evidence from Double Machine Learning
%% ============================================================
\documentclass[11pt]{article}

\usepackage[margin=1in, top=1.0in, bottom=1.0in]{geometry}
\usepackage{hyperref}
\usepackage{booktabs}
\usepackage{longtable}
\usepackage{array}
\usepackage{xcolor}
\usepackage{enumitem}
\usepackage{fancyhdr}
\usepackage{titlesec}
\usepackage{parskip}
\usepackage{listings}
\usepackage{microtype}
\usepackage{mdframed}
\usepackage{wasysym}

\hypersetup{
    colorlinks = true,
    linkcolor  = blue!65!black,
    urlcolor   = blue!65!black,
    citecolor  = blue!65!black,
    pdftitle   = {Reproducibility Package README},
}

\lstset{
    basicstyle       = \ttfamily\small,
    backgroundcolor  = \color{gray!7},
    frame            = single,
    rulecolor        = \color{gray!35},
    breaklines       = true,
    breakatwhitespace= true,
    columns          = flexible,
    keepspaces       = true,
    showstringspaces = false,
    tabsize          = 4,
}

\newmdenv[
    backgroundcolor  = yellow!8,
    linecolor        = orange!55,
    linewidth        = 1pt,
    innerleftmargin  = 8pt,
    innerrightmargin = 8pt,
    innertopmargin   = 6pt,
    innerbottommargin= 6pt,
    skipabove        = 5pt,
    skipbelow        = 5pt,
]{notebox}

\pagestyle{fancy}
\fancyhf{}
\lhead{\small Reproducibility Package README}
\rhead{\small Agricultural Loans \& DML --- Nigeria GHS}
\cfoot{\thepage}
\renewcommand{\headrulewidth}{0.4pt}

\titleformat{\section}{\large\bfseries}{}{0em}{}[\vspace{-4pt}\rule{\linewidth}{0.6pt}]
\titleformat{\subsection}{\normalsize\bfseries}{}{0em}{}

\setlist[itemize]{itemsep=1pt, topsep=3pt, leftmargin=*}
\setlist[enumerate]{itemsep=1pt, topsep=3pt, leftmargin=*}

%% ============================================================
\begin{document}

\begin{center}
    {\large\bfseries Agricultural Loans and Expenditure Allocation in Nigeria:\\[2pt]
    Evidence from Double Machine Learning} \\[8pt]
    {\normalsize [Author(s)]} \\[4pt]
    {\small \today}
\end{center}

\vspace{4pt}\hrule\vspace{8pt}

\tableofcontents
\newpage

%% ==============================================================
\section{Reproducibility Package README}
%% ==============================================================

The code within this replication package executes the analysis detailed in the
paper ``Agricultural Loans and Expenditure Allocation in Nigeria: Evidence from
Double Machine Learning'' using Stata, R, and Python. This package includes all
code required to go from the raw survey microdata to the replication of all
tables and figures in the analysis. The Stata master file (\texttt{00\_master.do})
manages sequential execution of all data construction code, anticipated to run
for approximately 30--45 minutes. Python estimation is run in a second stage.
Our primary objective is to ensure accessibility and credibility in replicating
our research methodology and results.

%% ==============================================================
\section{Data Availability and Provenance Statements}
%% ==============================================================

The analysis uses two rounds of Nigeria's General Household Survey (GHS), a
nationally representative panel survey conducted by the National Bureau of
Statistics (NBS) in collaboration with the World Bank. The paper draws on
Wave~4 (2018/19) and Wave~5 (2023/24). Raw GHS microdata are publicly
available from the World Bank Microdata Library
(\url{https://microdata.worldbank.org}); free registration is required.

In addition, the analysis uses pre-processed LSMS panel data compiled by the
Evans School Policy Analysis and Research Group (EPAR), available at:
\begin{center}
\url{https://github.com/EvansSchoolPolicyAnalysisAndResearch/LSMS-Data-Dissemination/tree/main/Nigeria\%20GHS}
\end{center}

\subsection{Statement about Rights}

I certify that the author(s) of the manuscript have legitimate access to and
permission to use the data used in this manuscript.

\subsection{Summary of Availability}

\begin{itemize}[label=\Square]
    \item[\CheckedBox] All data are publicly available.
    \item Some data cannot be made publicly available.
    \item No data can be made publicly available.
\end{itemize}

Both the World Bank GHS microdata and the EPAR processed files are freely and
publicly accessible at the links above. Raw data files (\texttt{.dta},
\texttt{.csv}) are excluded from this repository via \texttt{.gitignore}.
Place Wave~4 files in \texttt{Source Data/Wave 4/} and Wave~5 files in
\texttt{Source Data/Wave 5/} before running any code.

%% ==============================================================
\section{Getting Started}
%% ==============================================================

\subsection{Outline and Folder Structure}

The reproducibility package consists of five top-level folders:

\begin{itemize}
    \item \textbf{\texttt{Source Data/}} --- Raw GHS survey files
          (not tracked by the repository; download from the links above).
    \item \textbf{\texttt{Stata Code/}} --- All Stata do-files for data
          cleaning and table production. Final analysis dataset is written
          to \texttt{Stata Code/Stata Data Landing/DML Cleaned Data.dta}.
    \item \textbf{\texttt{R Code/}} --- Single script
          (\texttt{FIES Index Calculation.r}) for Rasch-model food
          insecurity index; called automatically by the Stata master.
    \item \textbf{\texttt{Python Code/}} --- DML estimation, sensitivity
          analysis, and figure scripts. Outputs are written to
          \texttt{Python Code/Python Table Landing/} and
          \texttt{Python Code/Python Figure Landing/}.
    \item \textbf{\texttt{Tables and Figures/}} --- Final paper outputs
          (\texttt{.tex} tables, \texttt{.png} figures).
\end{itemize}

\begin{lstlisting}
$ROOT/
|-- 00_master.do                  # Stata master (Stage 1 entry point)
|-- 00_figures.py                 # Python figure master (Stage 2)
|-- Source Data/Wave 4/           # GHS Wave 4 raw files
|-- Source Data/Wave 5/           # GHS Wave 5 raw files
|-- Stata Code/
|   |-- Stata Data Collection and Cleaning Scripts/
|   |-- Stata Control Covariate Collection and Cleaning Scripts/
|   |-- Stata Expenditure Calculations Scripts/
|   |-- Stata Output Tables Scripts/
|   `-- Stata Data Landing/DML Cleaned Data.dta  # Stage 1 output
|-- R Code/FIES Index Calculation.r
|-- Python Code/
|   |-- Python DML Scripts/       # fold_generator.py, learners.py,
|   |                             # helper_functions.py,
|   |                             # nuisance_function_residual_est.py,
|   |                             # DML Identification.py,
|   |                             # DML Benchmarking.py
|   |-- Pthon Figures Scripts/    # DML Figure A1--A3, Benchmark, Summary
|   |-- Python Table Landing/nuisance_cache/
|   `-- Python Figure Landing/
`-- Tables and Figures/
\end{lstlisting}

\subsection{How to Reproduce}

\medskip
\textbf{PART 1: Stata --- Data Construction.}
Open \texttt{00\_master.do}, set the \texttt{global root} macro at line~1
to the repository root path, then run the file in full:

\begin{lstlisting}
global root "C:/path/to/2nd-Year-Paper"
do 00_master.do
\end{lstlisting}

This installs required Stata packages, calls the R FIES script, and
produces \texttt{Stata Code/Stata Data Landing/DML Cleaned Data.dta}.
All subsequent stages read from that file only. Estimated runtime: 30--45 min.

\medskip
\textbf{PART 2: Python --- Install Dependencies.}
Install all required packages with a single command (see
Section~\ref{sec:python-install} for the full \texttt{requirements.txt}):

\begin{lstlisting}
pip install -r requirements.txt
\end{lstlisting}

\medskip
\textbf{PART 3: Python --- Estimation and Output.}
Run the following scripts in order. No path changes are needed; each script
resolves paths relative to its own file location.

\begin{itemize}
    \item \textbf{PART 3.1 --- Main DML results.}
          \texttt{Python Code/Python DML Scripts/DML Identification.py} \\
          Estimates ATEs, HTEs (loan $\times$ farm size), and GATEs by
          farm-size quartile for each outcome. Output is a timestamped
          \texttt{.txt} file in \texttt{Python Table Landing/}.
          Set \texttt{N\_REP=5} for a fast check; paper results use
          \texttt{N\_REP=30}. Runtime: $\approx$20 min (\texttt{N\_REP=5});
          2--6 hr (\texttt{N\_REP=30}, 4 cores).

    \item \textbf{PART 3.2 --- Sensitivity analysis.}
          \texttt{Python Code/Python DML Scripts/DML Benchmarking.py} \\
          Computes Rotnitzky--Vansteelandt omitted-variable bounds by
          re-estimating with each control dropped in turn.
          Output CSV files written to \texttt{Python Table Landing/}.
          Runtime: 30--90 min.

    \item \textbf{PART 3.3 --- Descriptive figures (A1--A3).}
          \texttt{00\_figures.py} (run from \texttt{\$ROOT}) \\
          Produces kernel density figures. Saved to
          \texttt{Python Figure Landing/}.

    \item \textbf{PART 3.4 --- Sensitivity forest plot.}
          \texttt{Python Code/Pthon Figures Scripts/DML Figure Benchmark.py} \\
          Reads benchmark CSVs from PART~3.2 and saves the forest plot.

    \item \textbf{PART 3.5 --- Summary statistics.}
          \texttt{Python Code/Pthon Figures Scripts/DML Summary Statistics.py} \\
          Produces farm-size quartile summary table; saved as
          \texttt{Python Figure Landing/quartile\_summary\_anova.csv}.
\end{itemize}

\begin{notebox}
\textbf{Cache behavior.} After PART~3.1 runs once, nuisance predictions are
cached in \texttt{Python Table Landing/nuisance\_cache/}. Subsequent runs
load the cache automatically (\texttt{re\_estimate\_nuisance=False} by
default), which is much faster. Set \texttt{re\_estimate\_nuisance=True}
only if you change the dataset, learners, or fold/rep parameters.
\end{notebox}

%% ==============================================================
\section{Dataset List}
%% ==============================================================

\begin{longtable}{@{}p{0.36\textwidth}p{0.40\textwidth}p{0.13\textwidth}@{}}
\toprule
Dataset & Description & Provided \\
\midrule
\endfirsthead
\toprule
Dataset & Description & Provided \\
\midrule
\endhead
\bottomrule
\endfoot
\texttt{Source Data/Wave 4/*.dta} &
    Raw GHS Wave~4 (2018/19) survey modules from the World Bank
    Microdata Library. &
    No \\[3pt]
\texttt{Source Data/Wave 5/*.dta} &
    Raw GHS Wave~5 (2023/24) survey modules. &
    No \\[3pt]
EPAR LSMS processed files &
    Pre-processed Nigeria GHS panel data from the EPAR LSMS
    Dissemination repository (linked above). &
    No \\[3pt]
\texttt{Stata Data Landing/\newline DML Cleaned Data.dta} &
    Analysis-ready balanced panel produced by Stage~1. Sole data
    input to all Python scripts. &
    No\\[-1pt]
    & \textit{(built by Stage~1)} & \\[3pt]
\texttt{nuisance\_cache/\newline loan\_hat\_K5\_R30*.npy} &
    Cached loan propensity score predictions,
    shape \texttt{(n\_obs, N\_REP)}. &
    Yes \\[3pt]
\texttt{nuisance\_cache/\newline outcome\_hat\_*\_K5\_R*.npy} &
    Cached outcome conditional mean predictions per outcome variable. &
    Yes \\
\end{longtable}

%% ==============================================================
\section{List of Tables and Programs}
%% ==============================================================

\begin{longtable}{@{}p{0.26\textwidth}p{0.36\textwidth}p{0.27\textwidth}@{}}
\toprule
Table / Figure & Program & Output file \\
\midrule
\endfirsthead
\toprule
Table / Figure & Program & Output file \\
\midrule
\endhead
\bottomrule
\endfoot
Balance Table &
    \texttt{Stata Output Tables Scripts/\newline Balance Table.do} &
    \texttt{Tables and Figures/\newline balance\_table.tex} \\[3pt]
Main DML Results (ATE, HTE, GATE) &
    \texttt{Python DML Scripts/\newline DML Identification.py} &
    \texttt{Python Table Landing/\newline DML Identification\_v*\_*.txt} \\[3pt]
Sensitivity Bounds &
    \texttt{Python DML Scripts/\newline DML Benchmarking.py} &
    \texttt{Python Table Landing/\newline benchmark\_*.csv} \\[3pt]
Summary Statistics &
    \texttt{Pthon Figures Scripts/\newline DML Summary Statistics.py} &
    \texttt{Python Figure Landing/\newline quartile\_summary\_anova.csv} \\[3pt]
Figure A1 (FIES KDE) &
    \texttt{Pthon Figures Scripts/\newline DML Figure A1.py} &
    \texttt{Python Figure Landing/\newline Figure\_A1.png} \\[3pt]
Figure A2 (Farm-size KDE) &
    \texttt{Pthon Figures Scripts/\newline DML Figure A2.py} &
    \texttt{Python Figure Landing/\newline Figure\_A2.png} \\[3pt]
Figure A3 (KDE by loan formality) &
    \texttt{Pthon Figures Scripts/\newline DML Figure A3.py} &
    \texttt{Python Figure Landing/\newline Figure\_A3.png} \\[3pt]
Sensitivity Forest Plot &
    \texttt{Pthon Figures Scripts/\newline DML Figure Benchmark.py} &
    \texttt{Python Figure Landing/\newline Figure\_Benchmark.png} \\
\end{longtable}

%% ==============================================================
\section{Computational Requirements}
\label{sec:requirements}
%% ==============================================================

\subsection{Software Requirements}

\subsubsection*{Stata}

Stata/MP or Stata/SE version~17 or later. The following user-written
packages are \emph{automatically installed} by \texttt{00\_master.do}
on the first run:

\begin{table}[h!]
\centering
\begin{tabular}{@{}lll@{}}
\toprule
Package & Source & Purpose \\
\midrule
\texttt{univar}    & SSC & Univariate statistics \\
\texttt{ietoolkit} & SSC & IE data management utilities \\
\texttt{winsor2}   & SSC & Winsorizing continuous variables \\
\bottomrule
\end{tabular}
\end{table}

\subsubsection*{R}

R version~4.0 or later. Called via shell from \texttt{00\_master.do} for
the FIES index computation. R must be on the system \texttt{PATH} or its
path must be set in the shell command in \texttt{00\_master.do}. Required
package: \texttt{RM.weights} (install via
\texttt{install.packages("RM.weights")}).

\subsubsection*{Python}

Python~3.9 or later (3.12 used for paper results). Table~\ref{tab:python-packages}
lists all required packages with the exact versions used.

\begin{table}[h!]
\centering
\caption{Required Python packages}
\label{tab:python-packages}
\begin{tabular}{@{}lll@{}}
\toprule
Package & Version & Purpose \\
\midrule
\texttt{numpy}        & 1.26.4  & Numerical arrays \\
\texttt{pandas}       & 2.2.2   & Data manipulation; reads \texttt{.dta} via \texttt{read\_stata} \\
\texttt{scikit-learn} & 1.7.2   & Ensemble learners and cross-validation \\
\texttt{doubleml}     & 0.11.1  & DML Partial Linear Regression \\
\texttt{statsmodels}  & 0.14.6  & Clustered OLS for GATE estimation \\
\texttt{matplotlib}   & 3.9.0   & All figures \\
\texttt{scipy}        & 1.16.2  & Kernel density estimation \\
\texttt{joblib}       & 1.5.2   & Parallel nuisance estimation \\
\texttt{pyarrow}      & 22.0.0  & Parquet I/O for nuisance cache \\
\bottomrule
\end{tabular}
\end{table}

\subsection{Installing Python Dependencies}
\label{sec:python-install}

Save the following as \texttt{requirements.txt} in the repository root,
then run the installation command below.

\textbf{\texttt{requirements.txt}:}
\begin{lstlisting}
numpy==1.26.4
pandas==2.2.2
scikit-learn==1.7.2
doubleml==0.11.1
statsmodels==0.14.6
matplotlib==3.9.0
scipy==1.16.2
joblib==1.5.2
pyarrow==22.0.0
\end{lstlisting}

\textbf{Installation (run once from the repository root):}
\begin{lstlisting}
pip install -r requirements.txt
\end{lstlisting}

To install without pinned versions (use if the above causes conflicts):
\begin{lstlisting}
pip install numpy pandas scikit-learn doubleml statsmodels matplotlib scipy joblib pyarrow
\end{lstlisting}

\begin{notebox}
\textbf{Virtual environment (recommended):}
\begin{lstlisting}
python -m venv dml_env
dml_env\Scripts\activate        # Windows
# source dml_env/bin/activate   # macOS / Linux
pip install -r requirements.txt
\end{lstlisting}
\end{notebox}

\subsection{Key Configuration Parameters}

Before running Stage~2, review these settings in
\texttt{Python DML Scripts/DML Identification.py}:

\begin{table}[h!]
\centering
\begin{tabular}{@{}llp{0.46\textwidth}@{}}
\toprule
Parameter & Paper value & Notes \\
\midrule
\texttt{N\_REP}               & 30           & Set to 5 for a fast verification run \\
\texttt{k\_fold\_vector}      & \texttt{[5]} & Cross-fitting folds \\
\texttt{N\_WORKERS}           & 4            & Set to available CPU cores \\
\texttt{re\_estimate\_nuisance} & \texttt{False} & \texttt{True} forces cache rebuild \\
\texttt{HTE}                  & \texttt{True}  & Loan $\times$ farm-size interaction \\
\texttt{HTE\_GATE}            & \texttt{True}  & Farm-size quartile GATEs \\
\bottomrule
\end{tabular}
\end{table}

\texttt{N\_REP} and \texttt{k\_fold\_vector} must be set to the same values
in \texttt{DML Benchmarking.py}. The paper was run on Windows~11 with 4
physical cores and 16~GB RAM.

%% ==============================================================
\section{References}
%% ==============================================================

Lars Vilhuber, Connolly, M., Koren, M., Llull, J., \& Morrow, P. (2022).
\textit{A template README for social science replication packages} (v1.1).
Zenodo. \url{https://doi.org/10.5281/zenodo.7293838}

Chernozhukov, V., Chetverikov, D., Demirer, M., Duflo, E., Hansen, C.,
Newey, W., \& Robins, J. (2018). Double/debiased machine learning for
treatment and structural parameters.
\textit{The Econometrics Journal}, 21(1), C1--C68.

Bach, P., Chernozhukov, V., Kurz, M.\ S., \& Spindler, M. (2022).
DoubleML --- An object-oriented implementation of double machine learning
in Python. \textit{Journal of Machine Learning Research}, 23(53), 1--6.

Evans School Policy Analysis and Research Group (EPAR). (2024).
\textit{LSMS Data Dissemination: Nigeria GHS}.
\url{https://github.com/EvansSchoolPolicyAnalysisAndResearch/LSMS-Data-Dissemination}

\end{document}
